
%(BEGIN_QUESTION)
% Copyright 2006, Tony R. Kuphaldt, released under the Creative Commons Attribution License (v 1.0)
% This means you may do almost anything with this work of mine, so long as you give me proper credit

Hva er den matematiske sammenhengen mellom ventilstørrelse (i tommer) og maksimal strømningskapasitet ($C_v$)? For eksempel: hvordan sammenlignes strømningskapasiteten til en 8-tommers ventil med en 4-tommers ventil, når alle andre faktorer er like?

\underbar{file i01372}
%(END_QUESTION)





%(BEGIN_ANSWER)

$C_v$ er omtrent proporsjonal med {\it kvadratet} av nominell ventilstørrelse. Tanken her er at tverrsnittsarealet strømningen går gjennom øker med kvadratet av strømningsdiameteren ($A = \pi r^2$).

Dette ser vi i formelen for relativ strømningskapasitet ($C_d$), som er omtrent konstant for en gitt ventiltype (kule kontra globe kontra spjeld osv.):

$$C_d = {C_v \over d^2}$$

Omskrevet for å relatere diameter ($d$) til strømningskoeffisient ($C_v$):

$$d^2 = {C_v \over C_d}$$

Ved konstant $C_d$-verdi vil strømningskoeffisienten ($C_v$) øke med {\it kvadratet} av økende diameter.


%(END_ANSWER)





%(BEGIN_NOTES)

%INDEX% Final Control Elements, valve: sizing

%(END_NOTES)

